\scanpage{32}\mspage{30}
qui se factorisent par une des sous-catégories
pleines $\underline{\deleted{\unsure{\mathscr{L}|\mathcal{L}}}\mathrm{Ch}}_n(\unsure{M|\mathcal{M}})$. On peut dire aussi que
si on considère le système inductif des
cat. $\underline{\mathrm{Ch}}_n(M)$ comme un ind-objet
de Cat, \deleted{\unsure{disons}} qu'on va noter
\[
\underline{\underline{\mathrm{Ch}}}_{\deleted{\unsure{n|\infty}}}(M)
\qquad\text{ind. objet de Cat},
\]
on a
\[
\underline{\mathrm{Ch}}(M(X))\overset{\sim}{\simeq}
\underline{\mathrm{Hom}}\bigl(X,\underline{\underline{\mathrm{Ch}}}_{\deleted{\unsure{n|\infty}}}(M)\bigr)\quad .
\]
Si on en vient \deleted{\unsure{aux}} à $\unsure{\Sigma|Z},\Sigma_X$, On trouve
aussitôt qu'un chemin dans $M(X)$ est un
$\Sigma_X$-chemin, ssi \deleted{\unsure{tous ses objets}} \inserted{le foncteur} correspondant
$X\unsure{\mathord{-}|\to}\underline{\mathrm{Ch}}(M)$ se factorise par $\underline{\mathrm{Ch}}_\Sigma(M)$.

Ainsi on trouve
\[
\underline{\mathrm{Ch}}_{\Sigma_X}(M(X))\simeq\underline{\mathrm{Hom}}^!(X,\underline{\mathrm{Ch}}_\Sigma(M)),
\]
où le signe $!$ a le même sens que plus bas.
Si on introduit le ind-objet $\underline{\underline{\mathrm{Ch}}}_\Sigma(M)$
\unsure{défini|définie} par les $\underline{\mathrm{Ch}}_{\Sigma,n}(M)$, on trouve donc
aussi
\[
\underline{\mathrm{Ch}}_{\Sigma_X}(M(X))\xrightarrow{\approx}
\underline{\mathrm{Hom}}\bigl(X,\underline{\underline{\mathrm{Ch}}}_\Sigma(M)\bigr)\quad .
\]
% archive sheet number: 31
